%% 1-19 かたち：三角・四角・円（10ページ、p.192〜201）
%% 第1巻のきまり：15 pt・行送り28 pt、分かち書き（文節の間に半角の空白）、漢字はすべてルビ、式は全角の＋＝で書く

\makeatletter
% 直角のしるし  \bkqrright{(かど)}{1つめの辺の向き（度）}{大きさ}
\newcommand\bkqrright[3]{\begin{scope}[shift={#1}, rotate=#2]
  \draw[bkTheme, line width=0.9pt] (#3,0) -- (#3,#3) -- (0,#3);\end{scope}}
% 三角じょうぎ  \bkqrjougi{(直角のかど)}{1:45°の形 / 2:30°と60°の形}{横の辺の長さ}{回す角度}
%   直角のかどから、横の辺は右へ、たての辺は上へのびる
\newcommand\bkqrjougi[4]{\begin{scope}[shift={#1}, rotate=#4]
  \pgfmathsetmacro\bkqr@a{#3}%
  \ifnum#2=1 \pgfmathsetmacro\bkqr@b{#3}\else\pgfmathsetmacro\bkqr@b{#3*0.5774}\fi
  \pgfmathsetmacro\bkqr@r{(\bkqr@a+\bkqr@b-sqrt(\bkqr@a*\bkqr@a+\bkqr@b*\bkqr@b))/2}%
  \pgfmathsetmacro\bkqr@k{0.42}%
  \fill[bkBlue!14, even odd rule] (0,0) -- (\bkqr@a,0) -- (0,\bkqr@b) -- cycle
    ({\bkqr@r*(1-\bkqr@k)},{\bkqr@r*(1-\bkqr@k)}) -- ({\bkqr@r+(\bkqr@a-\bkqr@r)*\bkqr@k},{\bkqr@r*(1-\bkqr@k)})
    -- ({\bkqr@r*(1-\bkqr@k)},{\bkqr@r+(\bkqr@b-\bkqr@r)*\bkqr@k}) -- cycle;
  \draw[bkBlue!70!black, line width=0.9pt, line join=round] (0,0) -- (\bkqr@a,0) -- (0,\bkqr@b) -- cycle;
  \draw[bkBlue!55!black, line width=0.6pt, line join=round]
    ({\bkqr@r*(1-\bkqr@k)},{\bkqr@r*(1-\bkqr@k)}) -- ({\bkqr@r+(\bkqr@a-\bkqr@r)*\bkqr@k},{\bkqr@r*(1-\bkqr@k)})
    -- ({\bkqr@r*(1-\bkqr@k)},{\bkqr@r+(\bkqr@b-\bkqr@r)*\bkqr@k}) -- cycle;
  \pgfmathtruncatemacro\bkqr@n{(\bkqr@a-6)/2}%
  \foreach \i in {1,...,\bkqr@n} {\pgfmathtruncatemacro\bkqr@t{mod(\i,5)==0}%
    \ifnum\bkqr@t=1 \draw[bkBlue!60!black, line width=0.4pt] (2*\i,0) -- (2*\i,1.6);
    \else\draw[bkBlue!60!black, line width=0.3pt] (2*\i,0) -- (2*\i,0.9);\fi}
\end{scope}}
% コンパス  \bkqrcompass{(はりの先)}{ひらき}{高さ}
\newcommand\bkqrcompass[3]{\begin{scope}[shift={#1}]
  \coordinate (bkqr@top) at (#2/2,#3);
  \draw[bkGray!70!black, line width=2.4pt, line cap=round] ($(0,0)!0.12!(bkqr@top)$) -- (bkqr@top);
  \draw[bkInk, line width=0.8pt] (0,0) -- ($(0,0)!0.14!(bkqr@top)$);
  \draw[bkGray!70!black, line width=2.4pt, line cap=round] ($(#2,0)!0.3!(bkqr@top)$) -- (bkqr@top);
  \draw[bkYellow!90!black, line width=2.6pt] ($(#2,0)!0.07!(bkqr@top)$) -- ($(#2,0)!0.3!(bkqr@top)$);
  \fill[bkInk] (#2,0) -- ($(#2,0)!0.08!(bkqr@top)+(-0.9,0)$) -- ($(#2,0)!0.08!(bkqr@top)+(0.9,0)$) -- cycle;
  \fill[bkGray!60, rounded corners=0.6mm] ($(bkqr@top)+(-1.5,0)$) rectangle ($(bkqr@top)+(1.5,6)$);
  \fill[bkRed] (bkqr@top) circle (1.6);
\end{scope}}
% 車（タイヤの形をえらべる）  \bkqrcar{(x,y)}{大きさ}{色}{0:まるいタイヤ / 1:四角いタイヤ}   (x,y) は左のタイヤの下
\newcommand\bkqrcar[4]{\begin{scope}[shift={#1}, scale=#2]
  \fill[#3, rounded corners=2.5] (-3,3.6) rectangle (27,10.6);
  \fill[#3, rounded corners=2.5] (3,9.6) -- (7,16.6) -- (19,16.6) -- (23,9.6) -- cycle;
  \fill[white, rounded corners=1] (8,10.8) -- (10.4,15.2) -- (12.4,15.2) -- (12.4,10.8) -- cycle (13.6,10.8) -- (13.6,15.2) -- (17.6,15.2) -- (21,10.8) -- cycle;
  \fill[bkYellow] (24.5,6.6) rectangle (27,8.6);
  \ifnum#4=0 \fill[bkInk!85] (3,3.3) circle (3.3) (21,3.3) circle (3.3);
  \else \fill[bkInk!85] (-0.3,0) rectangle (6.3,6.6) (17.7,0) rectangle (24.3,6.6);\fi
  \fill[bkLine] (3,3.3) circle (1.2) (21,3.3) circle (1.2);
\end{scope}}
% 形につける名前の札  \bkqrtag{(x,y)}{A}
\newcommand\bkqrtag[2]{\node[circle, draw=bkInk!70, fill=white, line width=0.7pt, text=bkInk,
  font=\sffamily\bfseries\footnotesize, minimum size=5.2mm, inner sep=0pt] at #1 {#2};}
% はぐるま  \bkqrgear{(x,y)}{半径}{歯の数}{色}
\newcommand\bkqrgear[4]{\begin{scope}[shift={#1}]
  \pgfmathsetmacro\bkqr@s{360/#3}%
  \foreach \k in {1,...,#3} {\fill[#4, rotate=\k*\bkqr@s] (-1.1,#2-0.6) rectangle (1.1,#2+1.6);}
  \fill[#4] (0,0) circle (#2);
  \fill[white] (0,0) circle (#2*0.32);
  \draw[#4!60!black, line width=0.6pt] (0,0) circle (#2*0.32);
\end{scope}}
\makeatother

\bkchapter{1}{19}{かたち：\R{三角}{さん|かく}・\R{四角}{し|かく}・\R{円}{えん}}{\bklevel{1}}{かたち：さんかく・しかく・えん}
\bkmichishirube{1-8,1-12}{1-20,1-21,4-8}

\begin{goal}
\begin{itemize}[label={\color{bkTheme}●}, leftmargin=1.4\zw, itemsep=0pt, topsep=0pt]
\item \R{三角形}{さん|かく|けい}と \R{四角形}{し|かく|けい}を \R{見}{み}わけられる。
\item \R{直角}{ちょっ|かく}を \R{見}{み}つけて、 \R{長方形}{ちょう|ほう|けい}や \R{正方形}{せい|ほう|けい}が わかる。
\item コンパスで \R{円}{えん}を かいて、 \R{中心}{ちゅう|しん}と \R{半径}{はん|けい}が わかる。
\end{itemize}
\end{goal}

\begin{hatena}
\begin{bkfigure}
\begin{tikzpicture}[x=1mm, y=1mm]
  % A：まるい タイヤ
  \fill[white, rounded corners=3mm] (0,0) rectangle (59,56);
  \node[circle, fill=bkYellow, text=white, font=\sffamily\bfseries, minimum size=6.5mm, inner sep=0pt] at (6,50) {A};
  \node[anchor=west, font=\small, text=bkInk] at (11,50) {まるい タイヤ};
  \fill[bkGray!45] (3,6) rectangle (56,9.5);
  \foreach \x in {6,18,30,42} {\fill[white] (\x,7.5) rectangle ++(6,0.6);}
  \bkqrcar{(9,9.5)}{1.55}{bkRoute6}{0}
  \foreach \y in {19,24,29} {\draw[bkGray!70, line width=0.9pt, line cap=round] (2,\y) -- (5.5,\y);}
  % B：四角い タイヤ
  \begin{scope}[shift={(63,0)}]
  \fill[white, rounded corners=3mm] (0,0) rectangle (59,56);
  \node[circle, fill=bkYellow, text=white, font=\sffamily\bfseries, minimum size=6.5mm, inner sep=0pt] at (6,50) {B};
  \node[anchor=west, font=\small, text=bkInk] at (11,50) {\R{四角}{し|かく}い タイヤ};
  \fill[bkGray!45] (3,6) rectangle (56,9.5);
  \foreach \x in {6,18,30,42} {\fill[white] (\x,7.5) rectangle ++(6,0.6);}
  \bkqrcar{(9,9.5)}{1.55}{bkRoute2}{1}
  \node[font=\sffamily\bfseries\Large, text=bkYellow!80!black] at (52,40) {？};
  \end{scope}
\end{tikzpicture}
\end{bkfigure}
\noindent A は まるい タイヤ、 B は \R{四角}{し|かく}い タイヤの \R{車}{くるま}です。\\
B の \R{車}{くるま}が はしったら、 どう なるかな？\\
タイヤは、 どうして まるいのかな？
\end{hatena}

\clearpage
%% ---------------- p.193 やってみよう ----------------
\begin{yattemiyou}[かたちを さがそう]
\begin{enumerate}[label={\bkcircled[bkGreen]{\arabic*}}, leftmargin=2.2\zw, itemsep=2pt, topsep=0pt]
\item まる・さんかく・しかくの ものは、 それぞれ いくつ あるかな？
\end{enumerate}
\begin{bkfigure}
\begin{tikzpicture}[x=1mm, y=1mm]
  \fill[white, rounded corners=3mm] (0,0) rectangle (122,48);
  % 1だんめ
  \bkclock*{(13,35)}{8}{3}{0}
  \fill[bkStem!70] (30,29) rectangle (48,42); \fill[bkStem!45] (31.5,29) rectangle (34,42);      % 本
  \draw[bkStem!90!black, line width=0.7pt] (30,29) rectangle (48,42);
  \fill[bkAppleRed, line join=round] (57,42) -- (73,42) -- (65,28.2) -- cycle;                       % とまれ
  \draw[white, line width=1pt, line join=round] (59.4,40.6) -- (70.6,40.6) -- (65,30.9) -- cycle;
  \node[text=white, font=\sffamily\bfseries\fontsize{5.6}{6}\selectfont] at (65,36.6) {とまれ};
  \bkcookie{(85,35)}{7}
  \fill[bkRoute6!80] (98,28) rectangle (112,42);                                                     % おりがみ
  % 2だんめ
  \fill[bkGreen!60!black] (6,8) rectangle (30,22); \fill[bkStem!60] (6,6.6) rectangle (30,8);       % こくばん
  \draw[white, line width=0.6pt] (9,17) -- (15,17) (9,13) -- (20,13);
  \filldraw[fill=bkYellow!60, draw=bkStem, line width=0.8pt, line join=round] (38,7) -- (54,7) -- (38,21) -- cycle;  % サンドイッチ
  \draw[bkLeaf, line width=1.6pt] (40.5,8.4) -- (50.5,8.4);
  \filldraw[fill=bkRoute4!70, draw=bkRoute4!60!black, line width=0.7pt] (66,14) circle (6.5);      % ボタン
  \foreach \p in {(64.2,15.8),(67.8,15.8),(64.2,12.2),(67.8,12.2)} {\fill[white] \p circle (0.9);}
  \filldraw[fill=white, draw=bkPurple, line width=1pt] (80,7.5) rectangle (100,20.5);              % カード
  \fill[bkPurple!70] (84,16) rectangle (96,18);\fill[bkPurple!35] (84,11) rectangle (93,13);
  \fill[bkPurple!35] (84,9) rectangle (90,10);
\end{tikzpicture}
\end{bkfigure}
\noindent\hspace*{1\zw}{\small\color{bkOrange!58!black}ヒント：ゆびで さしながら かぞえると、 かぞえわすれが ないよ。}\par
\noindent\hspace*{1\zw}まる\enspace\bkkakikomi\enspace こ\hspace{1.5\zw}さんかく\enspace\bkkakikomi\enspace こ\hspace{1.5\zw}しかく\enspace\bkkakikomi\enspace こ
\begin{enumerate}[label={\bkcircled[bkGreen]{\arabic*}}, leftmargin=2.2\zw, itemsep=2pt, topsep=4pt, start=2]
\item ラップの つつと、 ティッシュの はこを、 つくえの \R{上}{うえ}で ころがして みよう。
\end{enumerate}
\begin{bkfigure}
\begin{tikzpicture}[x=1mm, y=1mm]
  \fill[white, rounded corners=3mm] (0,0) rectangle (122,32);
  % つつ（ななめうしろへ のびる）
  \draw[bkInk!60, line width=0.8pt] (3,5) -- (58,5);
  \filldraw[fill=bkRoute6!35, draw=bkRoute6!80!black, line width=1pt, line join=round]
    (10.4,19.2) -- (20.4,24.2) arc (116.6:-63.4:8) -- (17.6,5.8) -- cycle;
  \filldraw[fill=bkRoute6!20, draw=bkRoute6!80!black, line width=1pt] (14,13) circle (8);
  \filldraw[fill=white, draw=bkRoute6!80!black, line width=0.8pt] (14,13) circle (5.4);
  \draw[bk arrow, bkGreen] (34,9) -- (45,9);
  \draw[bkGreen, line width=1.1pt, -{Stealth[length=5pt]}] (6,22) arc (140:70:10);
  \node[font=\sffamily\bfseries\small, text=bkGreen!70!black] at (46,20) {ころころ};
  % はこ（ティッシュの はこ）
  \begin{scope}[shift={(64,0)}]
  \draw[bkInk!60, line width=0.8pt] (0,5) -- (56,5);
  \filldraw[fill=bkYellow!45, draw=bkStem, line width=1pt, line join=round] (20,5) -- (30,10) -- (30,24) -- (20,19) -- cycle;
  \filldraw[fill=bkYellow!20, draw=bkStem, line width=1pt, line join=round] (4,19) -- (14,24) -- (30,24) -- (20,19) -- cycle;
  \fill[white] (12.5,20.6) -- (19.5,22.4) -- (21.5,22.4) -- (14.5,20.6) -- cycle;
  \filldraw[fill=bkYellow!30, draw=bkStem, line width=1pt] (4,5) rectangle (20,19);
  \draw[bkGreen, line width=1.1pt, -{Stealth[length=5pt]}] (2,24) arc (140:80:9);
  \node[font=\sffamily\bfseries\small, text=bkGreen!70!black] at (45,20) {ごとん？};
  \end{scope}
\end{tikzpicture}
\end{bkfigure}
\noindent\hspace*{1\zw}よく ころがるのは\enspace\bkkakikomi[30mm]
\end{yattemiyou}

\bkfukidashi{mina}{つつは ころころ、 はこは ごとん ごとん。\\まるい \R{形}{かたち}には、 なにか ひみつが ありそう！}

\clearpage
%% ---------------- p.194〜195 数学の道具 ----------------
\bksection{doguu}{\R{数学}{すう|がく}の \R{道具}{どう|ぐ}}

\bksubsection{1}{\R{三角形}{さん|かく|けい}と \R{四角形}{し|かく|けい}}

\label{bk:1-19:sankaku}まっすぐな ストローを つないで、 まわりを かこんで みよう。

\begin{bkfigure}
\begin{tikzpicture}[x=1mm, y=1mm]
  % 三角形
  \coordinate (a1) at (6,4); \coordinate (a2) at (44,4); \coordinate (a3) at (20,21);
  \foreach \p/\q/\c in {a1/a2/bkRoute6, a2/a3/bkRoute2, a3/a1/bkRoute4} {\draw[\c, line width=3.4pt, line cap=round, shorten <=1.6mm, shorten >=1.6mm] (\p) -- (\q);}
  \foreach \p in {a1,a2,a3} {\fill[bkInk!80] (\p) circle (1.3);}
  \node[anchor=north, font=\sffamily\bfseries\small] at (25,0.5) {\R{辺}{へん} 3\R{本}{ぼん}、 ちょう\R{点}{てん} 3こ};
  % 四角形
  \coordinate (b1) at (74,4); \coordinate (b2) at (112,4); \coordinate (b3) at (118,19); \coordinate (b4) at (80,21.5);
  \foreach \p/\q/\c in {b1/b2/bkRoute6, b2/b3/bkRoute2, b3/b4/bkRoute4, b4/b1/bkRoute5} {\draw[\c, line width=3.4pt, line cap=round, shorten <=1.6mm, shorten >=1.6mm] (\p) -- (\q);}
  \foreach \p in {b1,b2,b3,b4} {\fill[bkInk!80] (\p) circle (1.3);}
  \node[anchor=north, font=\sffamily\bfseries\small] at (96,0.5) {\R{辺}{へん} 4\R{本}{ほん}、 ちょう\R{点}{てん} 4こ};
  % ラベル
  \node[font=\sffamily\bfseries\small, text=bkTheme] (h) at (47,17.5) {\R{辺}{へん}};
  \draw[bkTheme, line width=0.8pt, -{Stealth[length=4pt]}] (h.south west) -- (34,10.5);
  \node[font=\sffamily\bfseries\small, text=bkTheme] (c) at (63,23) {ちょう\R{点}{てん}};
  \draw[bkTheme, line width=0.8pt, -{Stealth[length=4pt]}] (c.east) -- (78,22);
\end{tikzpicture}
\end{bkfigure}

\noindent まっすぐな \R{線}{せん}を \R{直線}{ちょく|せん}と いいます。\\
3\R{本}{ぼん}の \R{直線}{ちょく|せん}で かこまれた \R{形}{かたち}を \R{三角形}{さん|かく|けい}、\\
4\R{本}{ほん}の \R{直線}{ちょく|せん}で かこまれた \R{形}{かたち}を \R{四角形}{し|かく|けい}と いいます。\bkidx{さんかくけい}{\R{三角形}{さん|かく|けい}}\bkidx{しかくけい}{\R{四角形}{し|かく|けい}}\bkidx{ちょくせん}{\R{直線}{ちょく|せん}}\\
まわりの \R{直線}{ちょく|せん}を \R{辺}{へん}、 かどの \R{点}{てん}を ちょう\R{点}{てん}と いいます。\bkidx{へん}{\R{辺}{へん}}\bkidx{ちょうてん}{ちょう\R{点}{てん}}

\par\vspace{3pt}
\noindent 「かこまれた」とは、 まわりが ぐるっと つながって いて、 すきまが ない ことです。
「\R{三}{さん}」は 3、 「\R{四}{し}」は 4の ことです。
\R{三角形}{さん|かく|けい}は \R{辺}{へん}が 3\R{本}{ぼん}、 ちょう\R{点}{てん}が 3こ。
\R{四角形}{し|かく|けい}は \R{辺}{へん}が 4\R{本}{ほん}、 ちょう\R{点}{てん}が 4こです。
かぞえる ときは、 1つの ちょう\R{点}{てん}から \R{辺}{へん}を ゆびで なぞって、 もとの ちょう\R{点}{てん}に もどるまで かぞえよう。

\bksubsection{2}{\R{直角}{ちょっ|かく}}

\label{bk:1-19:chokkaku}\R{紙}{かみ}を 2\R{回}{かい} おろう。 2\R{回}{かい}\R{目}{め}は、 おりめを かさねて おるよ。

\begin{bkfigure}
\begin{tikzpicture}[x=0.8mm, y=0.8mm]
  % ①
  \filldraw[fill=bkYellow!35, draw=bkYellow!70!black, line width=0.8pt, line join=round]
    (0,4) -- (9,-1) -- (24,0) -- (32,6) -- (30,17) -- (21,23) -- (8,21) -- (1,14) -- cycle;
  \draw[bkGray, line width=0.8pt, dash pattern=on 2pt off 1.5pt] (-2,10) -- (34,10);
  \node[anchor=north, font=\sffamily\bfseries\footnotesize] at (16,-3) {\bkcircled[bkTheme]{1}\ 1\R{回}{かい} おる};
  \draw[bk arrow, bkGray] (37,10) -- (45,10);
  % ②
  \begin{scope}[shift={(48,0)}]
  \filldraw[fill=bkYellow!35, draw=bkYellow!70!black, line width=0.8pt, line join=round]
    (0,2) -- (33,2) -- (30,13) -- (21,19) -- (8,18) -- (1,12) -- cycle;
  \draw[bkYellow!40!black, line width=1.4pt] (0,2) -- (33,2);
  \draw[bkGray, line width=0.8pt, dash pattern=on 2pt off 1.5pt] (16,-1) -- (16,22);
  \node[anchor=north, font=\sffamily\bfseries\footnotesize, align=center] at (16,-3) {\bkcircled[bkTheme]{2}\ かさねて おる};
  \end{scope}
  \draw[bk arrow, bkGray] (85,10) -- (93,10);
  % ③
  \begin{scope}[shift={(98,0)}]
  \filldraw[fill=bkYellow!35, draw=bkYellow!70!black, line width=0.8pt, line join=round]
    (0,2) -- (18,2) -- (16,12) -- (9,19) -- (0,19) -- cycle;
  \draw[bkYellow!40!black, line width=1.4pt] (18,2) -- (0,2) -- (0,19);
  \bkqrright{(0,2)}{0}{3.5}
  \node[anchor=north, font=\sffamily\bfseries\footnotesize] at (9,-3) {\bkcircled[bkTheme]{3}\ できた かど};
  \end{scope}
\end{tikzpicture}
\end{bkfigure}

\noindent おって できた かどの \R{形}{かたち}を \R{直角}{ちょっ|かく}と いいます。\bkidx{ちょっかく}{\R{直角}{ちょっ|かく}}

\par\vspace{3pt}
\noindent\begin{minipage}[c]{0.58\linewidth}
1\R{回}{かい}\R{目}{め}の おりめは、 まっすぐな \R{直線}{ちょく|せん}です。
2\R{回}{かい}\R{目}{め}は、 その \R{直線}{ちょく|せん}を \R{半分}{はん|ぶん}に かさねて おりました。
だから できた かどは、 まっすぐを ちょうど \R{半分}{はん|ぶん}に した かどです。
\R{紙}{かみ}を ひらくと、 \R{直角}{ちょっ|かく}が 2つ ならんで、 まっすぐに もどります。
\end{minipage}\hfill
\begin{minipage}[c]{0.4\linewidth}\centering
\begin{tikzpicture}[x=1mm, y=1mm]
  \fill[bkYellow!25] (-20,0) rectangle (20,17);
  \draw[bkYellow!40!black, line width=1.4pt] (-20,0) -- (20,0);
  \draw[bkGray, line width=0.9pt, dash pattern=on 2pt off 1.5pt] (0,0) -- (0,17);
  \bkqrright{(0,0)}{0}{3.2}\bkqrright{(0,0)}{90}{3.2}
  \node[font=\sffamily\bfseries\fontsize{8}{9}\selectfont, text=bkTheme] at (-10,8) {\R{直角}{ちょっ|かく}};
  \node[font=\sffamily\bfseries\fontsize{8}{9}\selectfont, text=bkTheme] at (10,8) {\R{直角}{ちょっ|かく}};
  \node[anchor=north, font=\sffamily\fontsize{8}{9}\selectfont, text=bkInk] at (0,-1.2) {ひらいた \R{紙}{かみ}};
\end{tikzpicture}
\end{minipage}

\noindent\begin{minipage}[c]{0.5\linewidth}
\R{三角}{さん|かく}じょうぎにも、\\
\R{直角}{ちょっ|かく}の かどが あります。
\end{minipage}\hfill
\begin{minipage}[c]{0.48\linewidth}\raggedleft
\begin{tikzpicture}[x=1mm, y=1mm]
  \bkqrjougi{(0,0)}{1}{14}{0}\bkqrright{(0,0)}{0}{3}
  \bkqrjougi{(18,0)}{2}{24}{0}\bkqrright{(18,0)}{0}{3}
  \node[anchor=east, font=\sffamily\footnotesize, text=bkTheme, inner sep=1pt, align=center] (m) at (-3,8) {\R{直角}{ちょっ|かく}の\\しるし};
  \draw[bkTheme, line width=0.7pt, -{Stealth[length=3.5pt]}] (m.south) .. controls (-9,1) and (-4,1.5) .. (-0.6,1.5);
\end{tikzpicture}
\end{minipage}

\par\vspace{4pt}
\noindent 「これは \R{直角}{ちょっ|かく}かな？」と おもったら、 \R{三角}{さん|かく}じょうぎの \R{直角}{ちょっ|かく}の かどを、 しらべたい かどに \R{合}{あ}わせます。
\R{下}{した}の \R{辺}{へん}に、 じょうぎの \R{辺}{へん}を ぴったり かさねるのが こつです。

\begin{bkfigure}
\begin{tikzpicture}[x=1mm, y=1mm]
  \foreach \i/\a/\t in {0/90/{ぴったり\\→ \R{直角}{ちょっ|かく}},1/112/{すきまが ある\\→ \R{直角}{ちょっ|かく}より \R{大}{おお}きい},2/68/{じょうぎが はみ\R{出}{だ}す\\→ \R{直角}{ちょっ|かく}より \R{小}{ちい}さい}} {
    \begin{scope}[shift={(\i*44+10,0)}]
      \bkqrjougi{(0,0)}{1}{16}{0}
      \draw[bkInk!85, line width=1.3pt, line cap=round] (26,0) -- (0,0) -- (\a:24);
      \fill[bkInk] (0,0) circle (0.8);
      \node[anchor=north, font=\sffamily\bfseries\fontsize{8.5}{11}\selectfont, text=bkInk, align=center, inner sep=0pt] at (12,-2) {\t};
    \end{scope}}
\end{tikzpicture}
\end{bkfigure}

\bksubsection{3}{\R{直角}{ちょっ|かく}の ある \R{形}{かたち}}

\label{bk:1-19:katachi}\R{三角}{さん|かく}じょうぎで かどを しらべて、 なかまに \R{分}{わ}けよう。\bkidx{ちょうほうけい}{\R{長方形}{ちょう|ほう|けい}}\bkidx{せいほうけい}{\R{正方形}{せい|ほう|けい}}\bkidx{ちょっかくさんかくけい}{\R{直角三角形}{ちょっ|かく|さん|かく|けい}}

\begin{bkfigure}
\begin{tikzpicture}[x=1mm, y=1mm]
  % 長方形
  \filldraw[fill=bkRoute6!15, draw=bkRoute6!80!black, line width=1.2pt] (2,0) rectangle (38,18);
  \foreach \px/\py/\a in {2/0/0,38/0/90,38/18/180,2/18/270} {\bkqrright{(\px,\py)}{\a}{3}}
  \bkqrjougi{(38,0)}{1}{11}{90}
  \node[font=\sffamily\bfseries, text=bkInk] at (20,-6) {\R{長方形}{ちょう|ほう|けい}};
  \node[anchor=north, font=\sffamily\footnotesize, text=bkInk, align=center] at (20,-10.5) {かどが みんな\\\R{直角}{ちょっ|かく}};
  % 正方形
  \filldraw[fill=bkRoute2!15, draw=bkRoute2!80!black, line width=1.2pt] (52,0) rectangle (70,18);
  \foreach \px/\py/\a in {52/0/0,70/0/90,70/18/180,52/18/270} {\bkqrright{(\px,\py)}{\a}{3}}
  \node[font=\sffamily\bfseries, text=bkInk] at (61,-6) {\R{正方形}{せい|ほう|けい}};
  \node[anchor=north, font=\sffamily\footnotesize, text=bkInk, align=center] at (61,-10.5) {かどが みんな \R{直角}{ちょっ|かく}。\\\R{辺}{へん}の \R{長}{なが}さも みんな \R{同}{おな}じ};
  % 直角三角形
  \filldraw[fill=bkRoute4!15, draw=bkRoute4!80!black, line width=1.2pt, line join=round] (88,0) -- (120,0) -- (88,18) -- cycle;
  \bkqrright{(88,0)}{0}{3}
  \node[font=\sffamily\bfseries, text=bkInk] at (104,-6) {\R{直角三角形}{ちょっ|かく|さん|かく|けい}};
  \node[anchor=north, font=\sffamily\footnotesize, text=bkInk, align=center] at (104,-10.5) {\R{直角}{ちょっ|かく}の かどが\\ある \R{三角形}{さん|かく|けい}};
\end{tikzpicture}
\end{bkfigure}

\noindent\begin{minipage}[c]{0.58\linewidth}
\R{方眼}{ほう|がん}に かいて \R{辺}{へん}の \R{長}{なが}さを しらべると、 \R{長方形}{ちょう|ほう|けい}は、 むかいあう \R{辺}{へん}の \R{長}{なが}さが \R{同}{おな}じです（\R{上}{うえ}と \R{下}{した}が 5マス、 \R{左}{ひだり}と \R{右}{みぎ}が 3マス）。
\R{正方形}{せい|ほう|けい}は、 4\R{本}{ほん}の \R{辺}{へん}が みんな \R{同}{おな}じ \R{長}{なが}さです。
\end{minipage}\hfill
\begin{minipage}[c]{0.4\linewidth}\centering
\begin{tikzpicture}[x=1mm, y=1mm]
  \bkgrid{(0,0)}{7}{5}{5}
  \filldraw[fill=bkRoute6!15, draw=bkRoute6!80!black, line width=1.2pt] (5,5) rectangle (30,20);
  \node[anchor=south, font=\sffamily\bfseries\fontsize{8}{9}\selectfont, fill=white, inner sep=0.6pt] at (17.5,20.6) {5マス};
  \node[anchor=north, font=\sffamily\bfseries\fontsize{8}{9}\selectfont, fill=white, inner sep=0.6pt] at (17.5,4.4) {5マス};
  \node[anchor=east, font=\sffamily\bfseries\fontsize{8}{9}\selectfont, inner sep=0.6pt] at (-0.6,12.5) {3マス};
  \node[anchor=west, font=\sffamily\bfseries\fontsize{8}{9}\selectfont, inner sep=0.6pt] at (35.6,12.5) {3マス};
\end{tikzpicture}
\end{minipage}

\Needspace{10\baselineskip}\bksubsection{4}{\R{円}{えん}と タイヤ}

\label{bk:1-19:en}\noindent\begin{minipage}[c]{0.64\linewidth}
コンパスで \R{円}{えん}を かこう。\bkidx{えん}{\R{円}{えん}}\\
はりの \R{点}{てん}が \R{中心}{ちゅう|しん}、 ひらきが \R{半径}{はん|けい}。\bkidx{はんけい}{\R{半径}{はん|けい}}\\
\R{半径}{はん|けい}は どこも \R{同}{おな}じ \R{長}{なが}さ。 コンパスの ひらきを かえずに まわすので、 はりから えんぴつまでの \R{長}{なが}さが いつも \R{同}{おな}じ だからです。\par
\R{中心}{ちゅう|しん}を とおって、 はしから はしまで ひいた \R{直線}{ちょく|せん}を \R{直径}{ちょっ|けい}と いいます。\bkidx{ちょっけい}{\R{直径}{ちょっ|けい}}\par
\R{直径}{ちょっ|けい}は、 \R{半径}{はん|けい}の 2つ\R{分}{ぶん}です。
\end{minipage}\hfill
\begin{minipage}[c]{0.34\linewidth}\centering
\begin{tikzpicture}[x=1mm, y=1mm]
  \draw[bkInk!80, line width=1.1pt] (0,0) circle (12);
  \draw[bkBlue, line width=1.3pt] (-12,0) -- (12,0);
  \foreach \xa/\xb in {-11.6/-0.7,0.7/11.6} {\draw[bkTheme, line width=0.6pt, decorate, decoration={brace, amplitude=1.1mm, mirror}] (\xa,-0.9) -- (\xb,-0.9)
    node[midway, below=1.6mm, font=\sffamily\bfseries\footnotesize, text=bkTheme, inner sep=0pt] {\R{半径}{はん|けい}};}
  \draw[bkBlue, line width=0.6pt, decorate, decoration={brace, amplitude=1.4mm, mirror}] (-12,-13.4) -- (12,-13.4)
    node[midway, below=2mm, font=\sffamily\bfseries\footnotesize, text=bkBlue, inner sep=0pt] {\R{直径}{ちょっ|けい}};
  \bkqrcompass{(0,0)}{12}{20}
  \fill[bkRed] (0,0) circle (1.1);
  \node[anchor=south east, font=\sffamily\bfseries\footnotesize, text=bkRed!80!black, inner sep=0.5pt] at (-1,1) {\R{中心}{ちゅう|しん}};
\end{tikzpicture}\par\vspace{2pt}
{\fontsize{9}{12}\selectfont\color{bkRed!80!black}\sffamily\bfseries ！\,はりは とがって いるよ。\\ \R{人}{ひと}に むけないでね。\par}
\end{minipage}

\begin{bkfigure}
\begin{tikzpicture}[x=0.82mm, y=0.82mm]
  % 円が ころがる
  \draw[bkInk!60, line width=0.9pt] (0,0) -- (60,0);
  \foreach \cx/\a in {8/90,30/-90,52/90} {
    \filldraw[fill=bkInk!12, draw=bkInk!80, line width=1.1pt] (\cx,7) circle (7);
    \draw[bkTheme, line width=1pt] (\cx,7) -- ++(\a:7);
    \fill[bkRed] (\cx,7) circle (1);}
  \draw[bkRed, line width=1pt, dash pattern=on 3pt off 1.5pt] (0,7) -- (60,7);
  \node[anchor=north, font=\sffamily\bfseries\footnotesize, text=bkInk] at (30,-1.5) {\R{中心}{ちゅう|しん}の \R{高}{たか}さは かわらない};
  % 四角が ころがる
  \begin{scope}[shift={(68,0)}]
  \draw[bkInk!60, line width=0.9pt] (0,0) -- (64,0);
  \draw[bkInk!25, line width=0.6pt, dash pattern=on 2pt off 2pt] (0,7) -- (64,7);
  \filldraw[fill=bkInk!12, draw=bkInk!80, line width=1.1pt] (4,0) rectangle (18,14);
  \filldraw[fill=bkInk!12, draw=bkInk!80, line width=1.1pt] (32,0) -- (41.9,9.9) -- (32,19.8) -- (22.1,9.9) -- cycle;
  \filldraw[fill=bkInk!12, draw=bkInk!80, line width=1.1pt] (46,0) rectangle (60,14);
  \draw[bkRed, line width=1pt] (11,7) arc (135:45:9.9) arc (135:45:9.9) arc (135:45:9.9);
  \foreach \p in {(11,7),(32,9.9),(53,7)} {\fill[bkRed] \p circle (1);}
  \node[anchor=north, font=\sffamily\bfseries\footnotesize, text=bkInk] at (32,-1.5) {\R{中心}{ちゅう|しん}が \R{上}{あ}がったり \R{下}{さ}がったり};
  \end{scope}
\end{tikzpicture}
\end{bkfigure}

\noindent \R{円}{えん}が ころがる とき、 \R{地面}{じ|めん}から \R{中心}{ちゅう|しん}までは、 いつも \R{半径}{はん|けい}の \R{長}{なが}さです。
\R{四角}{し|かく}は、 \R{辺}{へん}で \R{立}{た}つ ときと、 かどで \R{立}{た}つ ときで、 \R{中心}{ちゅう|しん}の \R{高}{たか}さが ちがいます。

\begin{bkcenter}
\begin{tblr}{colspec={Q[l,m,wd=50mm] Q[c,m,wd=30mm] Q[c,m,wd=30mm]}, hlines={0.6pt,bkGray!50}, vlines={0.6pt,bkGray!50}, rowsep=2pt, colsep=4pt, row{1}={bg=bkTheme!10, font=\sffamily\bfseries\small}, cells={font=\small}}
\R{中心}{ちゅう|しん}の \R{高}{たか}さ & \R{辺}{へん}・\R{円}{えん}で \R{立}{た}つ & かどで \R{立}{た}つ \\
\R{半径}{はん|けい} 5\,cm の \R{円}{えん} & 5\,cm & （かどが ない） \\
1\R{辺}{へん} 10\,cm の \R{正方形}{せい|ほう|けい} & 5\,cm & 7\,cm ほど \\
\end{tblr}
\end{bkcenter}

\noindent \R{四角}{し|かく}い タイヤは、 ころがる たびに \R{車}{くるま}が 5\,cm と 7\,cm ほどの \R{間}{あいだ}を \R{上}{あ}がったり \R{下}{さ}がったり します。
\bkhatenaref の \R{答}{こた}え：B の \R{車}{くるま}は、 がたがた ゆれる。

\begin{taisetsu}
\R{円}{えん}の \R{半径}{はん|けい}は、 どこも \R{同}{おな}じ \R{長}{なが}さ。\\
だから、 ころがしても \R{中心}{ちゅう|しん}の \R{高}{たか}さが かわらない。
\end{taisetsu}

\par\vspace{6pt}
%% ---------------- p.196〜197 例題 ----------------
\begin{reidai}[\R{三角形}{さん|かく|けい}と \R{四角形}{し|かく|けい}]
\R{三角形}{さん|かく|けい}は どれ？ \R{四角形}{し|かく|けい}は どれ？\par
\noindent\begin{minipage}[c]{0.48\linewidth}
\begin{tikzpicture}[x=1mm, y=1mm]
  \filldraw[fill=white, draw=bkInk!80, line width=1.1pt, line join=round] (0,0) -- (12,0) -- (3.5,13) -- cycle;
  \bkqrtag{(1,17)}{A}
  \filldraw[fill=white, draw=bkInk!80, line width=1.1pt, line join=round] (16,0) -- (28,1) -- (27,13) -- (17,11) -- cycle;
  \bkqrtag{(17,17)}{B}
  \filldraw[fill=white, draw=bkInk!80, line width=1.1pt, line join=round] (32,0) -- (44,0) .. controls (43,9) and (40,13) .. (36,13) -- cycle;
  \bkqrtag{(33,17)}{C}
  \draw[bkInk!80, line width=1.1pt, line join=round] (50.4,2.4) -- (48,0) -- (60,0) -- (54,13) -- (51.9,8);
  \bkqrtag{(49,17)}{D}
\end{tikzpicture}
\end{minipage}\hfill
\begin{minipage}[c]{0.5\linewidth}
\kangaekata \R{直線}{ちょく|せん}の\\\R{数}{かず}を かぞえる。
\end{minipage}\par
\kotae \R{三角形}{さん|かく|けい} A\quad \R{四角形}{し|かく|けい} B
\end{reidai}

\begin{jibun}
\R{三角形}{さん|かく|けい}は どれ？ \R{四角形}{し|かく|けい}は どれ？\par
\noindent\begin{minipage}[c]{0.58\linewidth}
\begin{tikzpicture}[x=1mm, y=1mm]
  \filldraw[fill=white, draw=bkInk!80, line width=1.1pt, line join=round] (0,1) -- (13,0) -- (14,11) -- (2,12) -- cycle;
  \bkqrtag{(1,17)}{A}
  \draw[bkInk!80, line width=1.1pt, line join=round] (23,12) -- (19,0) -- (33,0) -- (26,12.6);
  \bkqrtag{(20,17)}{B}
  \filldraw[fill=white, draw=bkInk!80, line width=1.1pt, line join=round] (38,0) -- (52,0) -- (48,13) -- cycle;
  \bkqrtag{(39,17)}{C}
  \filldraw[fill=white, draw=bkInk!80, line width=1.1pt] (65,6.5) circle (6.5);
  \bkqrtag{(59,17)}{D}
\end{tikzpicture}
\end{minipage}\hfill
\begin{minipage}[c]{0.4\linewidth}
\R{三角形}{さん|かく|けい}\enspace\bkkakikomi[16mm]\\
\R{四角形}{し|かく|けい}\enspace\bkkakikomi[16mm]
\end{minipage}
\bkjans{\mbox{\R{三角形}{さん|かく|けい} C}\quad \mbox{\R{四角形}{し|かく|けい} A}}
\end{jibun}

\begin{reidai}[\R{直角}{ちょっ|かく}を さがす]
\noindent\begin{minipage}[c]{0.56\linewidth}
\R{直角}{ちょっ|かく}は どれ？\par\vspace{1mm}
\begin{tikzpicture}[x=1mm, y=1mm]
  \bkqrjougi{(0,0)}{1}{10}{0}
  \draw[bkInk!80, line width=1.2pt, line join=round] (13,0) -- (0,0) -- (0,13);
  \bkqrtag{(4,16)}{A}
  \draw[bkInk!80, line width=1.2pt, line join=round] (32,0) -- (19,0) -- (25.5,11.3);
  \bkqrtag{(19,15)}{B}
  \draw[bkInk!80, line width=1.2pt, line join=round] (38,3.9) -- (44,0.5) -- (51,12.9);
  \bkqrtag{(39,15)}{C}
  \draw[bkInk!80, line width=1.2pt, line join=round] (67,0) -- (57,0) -- (54,11);
  \bkqrtag{(58,15)}{D}
\end{tikzpicture}
\end{minipage}\hfill
\begin{minipage}[c]{0.42\linewidth}
\kangaekata\\\R{三角}{さん|かく}じょうぎの\\\R{直角}{ちょっ|かく}を あてて みる。
\end{minipage}\par
\kotae A と C\quad （C は かたむいて いても \R{直角}{ちょっ|かく}）
\end{reidai}

\begin{jibun}
\noindent\begin{minipage}[c]{0.56\linewidth}
\begin{tikzpicture}[x=1mm, y=1mm]
  \draw[bkInk!80, line width=1.2pt, line join=round] (13,0) -- (0,0) -- (9.2,9.2);
  \bkqrtag{(2,13)}{A}
  \draw[bkInk!80, line width=1.2pt, line join=round] (17.5,9) -- (22,-1) -- (32,3.5);
  \bkqrtag{(17,15)}{B}
  \draw[bkInk!80, line width=1.2pt, line join=round] (51,0) -- (39,0) -- (36.5,12);
  \bkqrtag{(40,15)}{C}
  \draw[bkInk!80, line width=1.2pt, line join=round] (66,12) -- (66,0) -- (54,0);
  \bkqrtag{(60,14)}{D}
\end{tikzpicture}
\end{minipage}\hfill
\begin{minipage}[c]{0.42\linewidth}
\R{直角}{ちょっ|かく}は どれ？
\bkkotaeran[24mm]
\end{minipage}
\bkjans{B と D}
\end{jibun}

\begin{reidai}[\R{形}{かたち}の \R{名前}{な|まえ}]
\R{方眼}{ほう|がん}に かいた \R{形}{かたち}の \R{名前}{な|まえ}は？\par
\noindent\begin{minipage}[c]{0.52\linewidth}
\begin{tikzpicture}[x=1mm, y=1mm]
  \bkgrid{(0,0)}{13}{4}{5}
  \filldraw[fill=bkRoute6!20, draw=bkRoute6!80!black, line width=1.2pt] (5,5) rectangle (25,15);
  \bkqrtag{(2.5,17.5)}{A}
  \filldraw[fill=bkRoute2!20, draw=bkRoute2!80!black, line width=1.2pt] (30,5) rectangle (40,15);
  \bkqrtag{(27.5,17.5)}{B}
  \filldraw[fill=bkRoute4!20, draw=bkRoute4!80!black, line width=1.2pt, line join=round] (45,5) -- (60,5) -- (45,15) -- cycle;
  \bkqrtag{(47.5,17.5)}{C}
\end{tikzpicture}
\end{minipage}\hfill
\begin{minipage}[c]{0.46\linewidth}
\kangaekata かどと、\\\R{辺}{へん}の \R{長}{なが}さを しらべる。
\end{minipage}\par
\kotae A \R{長方形}{ちょう|ほう|けい}\quad B \R{正方形}{せい|ほう|けい}\quad C \R{直角三角形}{ちょっ|かく|さん|かく|けい}
\end{reidai}

\begin{jibun}
\noindent\begin{minipage}[c]{0.52\linewidth}
\begin{tikzpicture}[x=1mm, y=1mm]
  \bkgrid{(0,0)}{13}{4}{5}
  \filldraw[fill=bkRoute4!20, draw=bkRoute4!80!black, line width=1.2pt, line join=round] (5,15) -- (5,5) -- (15,5) -- cycle;
  \bkqrtag{(11.5,14.5)}{A}
  \filldraw[fill=bkRoute2!20, draw=bkRoute2!80!black, line width=1.2pt] (20,0) rectangle (35,15);
  \bkqrtag{(27.5,18)}{B}
  \filldraw[fill=bkRoute6!20, draw=bkRoute6!80!black, line width=1.2pt] (40,5) rectangle (60,10);
  \bkqrtag{(42.5,13.5)}{C}
\end{tikzpicture}
\end{minipage}\hfill
\begin{minipage}[c]{0.46\linewidth}
\R{形}{かたち}の \R{名前}{な|まえ}を かこう。
\end{minipage}\par
\noindent A\enspace\bkkakikomi[30mm]\hspace{0.8\zw}B\enspace\bkkakikomi[30mm]\hspace{0.8\zw}C\enspace\bkkakikomi[30mm]\par\vspace{4pt}
\noindent\begin{minipage}[c]{0.52\linewidth}
\begin{tikzpicture}[x=1mm, y=1mm]
  \bkgrid{(0,0)}{13}{4}{5}
  \filldraw[fill=bkRoute6!20, draw=bkRoute6!80!black, line width=1.2pt] (5,0) rectangle (15,20);
  \bkqrtag{(10,23.5)}{D}
  \filldraw[fill=bkRoute4!20, draw=bkRoute4!80!black, line width=1.2pt, line join=round] (20,5) -- (40,5) -- (40,15) -- cycle;
  \bkqrtag{(27,15)}{E}
  \filldraw[fill=bkRoute2!20, draw=bkRoute2!80!black, line width=1.2pt] (45,0) rectangle (65,20);
  \bkqrtag{(55,23.5)}{F}
\end{tikzpicture}
\end{minipage}\hfill
\begin{minipage}[c]{0.46\linewidth}
\R{方眼}{ほう|がん}の マスを かぞえて、 \R{辺}{へん}の \R{長}{なが}さも しらべよう。
\end{minipage}\par
\noindent D\enspace\bkkakikomi[30mm]\hspace{0.8\zw}E\enspace\bkkakikomi[30mm]\hspace{0.8\zw}F\enspace\bkkakikomi[30mm]
\bkjans{A \R{直角三角形}{ちょっ|かく|さん|かく|けい}、 B \R{正方形}{せい|ほう|けい}、 C \R{長方形}{ちょう|ほう|けい}、 D \R{長方形}{ちょう|ほう|けい}（たて 4マス、 よこ 2マス）、 E \R{直角三角形}{ちょっ|かく|さん|かく|けい}、 F \R{正方形}{せい|ほう|けい}（たて・よこ 4マス）}
\end{jibun}

\begin{reidai}[\R{形}{かたち}の \R{中}{なか}の \R{直角}{ちょっ|かく}を かぞえる]
この \R{形}{かたち}に、 \R{直角}{ちょっ|かく}は いくつ ありますか。\par
\noindent\begin{minipage}[c]{0.44\linewidth}\centering
\begin{tikzpicture}[x=1mm, y=1mm]
  \filldraw[fill=bkYellow!20, draw=bkInk!80, line width=1.2pt, line join=round] (0,0) -- (30,0) -- (30,16) -- (10,16) -- cycle;
  \bkqrright{(30,0)}{90}{3}\bkqrright{(30,16)}{180}{3}
  \bkqrtag{(-3.5,-1)}{A}\bkqrtag{(34,-1)}{B}\bkqrtag{(34,18)}{C}\bkqrtag{(7,19.5)}{D}
\end{tikzpicture}
\end{minipage}\hfill
\begin{minipage}[c]{0.54\linewidth}
\kangaekata ちょう\R{点}{てん}は 4こ。 A、 B、 C、 D の じゅんに、 \R{三角}{さん|かく}じょうぎを あてる。
\end{minipage}\par
\noindent A：はみ\R{出}{だ}す（\R{直角}{ちょっ|かく}より \R{小}{ちい}さい）。 B・C：ぴったり（\R{直角}{ちょっ|かく}）。 D：すきまが ある（\R{直角}{ちょっ|かく}より \R{大}{おお}きい）。\par
\kotae 2つ（B と C）
\end{reidai}

\begin{jibun}
\noindent\begin{minipage}[c]{0.44\linewidth}\centering
\begin{tikzpicture}[x=1mm, y=1mm]
  \filldraw[fill=bkRoute4!15, draw=bkInk!80, line width=1.2pt, line join=round] (0,0) -- (24,0) -- (24,12) -- (12,20) -- (0,12) -- cycle;
\end{tikzpicture}
\end{minipage}\hfill
\begin{minipage}[c]{0.54\linewidth}
この \R{形}{かたち}に、 \R{直角}{ちょっ|かく}は いくつ？ ちょう\R{点}{てん}を 1つずつ しらべよう。
\bkkotaeran[24mm]
\end{minipage}
\bkjans{2つ（\R{下}{した}の 2つの かど）}
\end{jibun}

\begin{reidai}[\R{円}{えん}の \R{半径}{はん|けい}と \R{直径}{ちょっ|けい}]
\R{半径}{はん|けい} 2\,cm の \R{円}{えん}が 3こ。 はしから はしまでは \R{何}{なん}cm？\par
\noindent\begin{minipage}[c]{0.5\linewidth}
\begin{tikzpicture}[x=1mm, y=1mm]
  \foreach \x in {9,27,45} {\draw[bkInk!80, line width=1.1pt] (\x,0) circle (9); \fill[bkRed] (\x,0) circle (0.9);}
  \draw[bkTheme, line width=1.1pt] (9,0) -- (18,0);
  \node[font=\sffamily\bfseries\fontsize{9.5}{10}\selectfont, text=bkTheme, inner sep=0pt] at (13.2,2.4) {2\,cm};
  \bkbrace{(0,-10.5)}{(54,-10.5)}{？\,cm}{-1}
\end{tikzpicture}
\end{minipage}\hfill
\begin{minipage}[c]{0.48\linewidth}
\kangaekata \R{直径}{ちょっ|けい}は \R{半径}{はん|けい}の 2つ\R{分}{ぶん}。\\2 × 2 ＝ 4で 4\,cm。\\\R{円}{えん}が すきまなく ならぶので、 \R{直径}{ちょっ|けい} 3つ\R{分}{ぶん}。
\end{minipage}\par
\shiki 4 × 3 ＝ 12 \kotaenext 12\,cm
\end{reidai}

\begin{jibun}
（1）\R{半径}{はん|けい} 4\,cm の \R{円}{えん}が 2こ。 はしから はしまでは \R{何}{なん}cm？
\bkshikiran
（2）\R{半径}{はん|けい} 3\,cm の \R{円}{えん}が 4こ。 はしから はしまでは \R{何}{なん}cm？
\bkshikiran
\bkjans{（1）8 × 2 ＝ 16\quad 16\,cm\quad \mbox{（\R{直径}{ちょっ|けい}は 8\,cm）}\quad（2）6 × 4 ＝ 24\quad 24\,cm\quad \mbox{（\R{直径}{ちょっ|けい}は 6\,cm）}}
\end{jibun}

\begin{reidai}[\R{直径}{ちょっ|けい}から \R{半径}{はん|けい}を もとめる]
\R{直径}{ちょっ|けい} 6\,cm の \R{円}{えん}を かきたい。 コンパスを \R{何}{なん}cm ひらけば よい？\par
\noindent\begin{minipage}[c]{0.4\linewidth}\centering
\begin{tikzpicture}[x=1mm, y=1mm]
  \draw[bkInk!80, line width=1.1pt] (0,0) circle (12);
  \draw[bkBlue, line width=1.2pt] (-12,0) -- (12,0);
  \fill[bkRed] (0,0) circle (0.9);
  \bkbrace{(-12,-14)}{(12,-14)}{6\,cm}{-1}
  \draw[bkTheme, line width=0.6pt, decorate, decoration={brace, amplitude=1.1mm}] (0.7,1) -- (11.6,1)
    node[midway, above=1.6mm, font=\sffamily\bfseries\footnotesize, text=bkTheme, inner sep=0pt] {？\,cm};
\end{tikzpicture}
\end{minipage}\hfill
\begin{minipage}[c]{0.58\linewidth}
\kangaekata コンパスの ひらきは \R{半径}{はん|けい}。\\
\R{直径}{ちょっ|けい}は \R{半径}{はん|けい}の 2つ\R{分}{ぶん}なので、 6を \R{同}{おな}じ 2つに \R{分}{わ}ける。\\
3 ＋ 3 ＝ 6（6 ÷ 2 ＝ 3、 1-14）。
\end{minipage}\par
\kotae 3\,cm
\end{reidai}

\begin{jibun}
\R{直径}{ちょっ|けい} 10\,cm の \R{円}{えん}を かくには、 コンパスを \R{何}{なん}cm ひらけば よい？
\bkkotaeran[24mm]
\bkjans{5\,cm（5 ＋ 5 ＝ 10。 10 ÷ 2 ＝ 5）}
\end{jibun}

\par\vspace{6pt}
%% ---------------- p.198 まちがえやすいところ ----------------
\bksection{tsumazuki}{まちがえやすい ところ}

\begin{tsumazuki}{すきまや まがった \R{線}{せん}}
\begin{minipage}[c]{0.62\linewidth}
すきまや まがった \R{線}{せん}が あると、\\\R{三角形}{さん|かく|けい}では ありません。
\end{minipage}\hfill
\begin{minipage}[c]{0.35\linewidth}\centering
\begin{tikzpicture}[x=1mm, y=1mm]
  \draw[bkInk!80, line width=1.1pt, line join=round] (3,3) -- (0,0) -- (16,0) -- (7,13) -- (4.5,8);
  \node at (8,-4.5) {\bkbatsu};
  \draw[bkInk!80, line width=1.1pt, line join=round] (22,0) -- (38,0) .. controls (37,9) and (32,13) .. (27,13) -- cycle;
  \node at (30,-4.5) {\bkbatsu};
\end{tikzpicture}
\end{minipage}
\end{tsumazuki}

\begin{tsumazuki}{かたむけても \R{正方形}{せい|ほう|けい}}
\begin{minipage}[c]{0.62\linewidth}
かたむけても、 \R{辺}{へん}の \R{長}{なが}さと \R{直角}{ちょっ|かく}は そのまま。 \R{正方形}{せい|ほう|けい}です。
\end{minipage}\hfill
\begin{minipage}[c]{0.35\linewidth}\centering
\begin{tikzpicture}[x=1mm, y=1mm]
  \filldraw[fill=bkRoute2!15, draw=bkRoute2!80!black, line width=1.1pt] (0,0) rectangle (11,11);
  \draw[bk arrow, bkGray, line width=1pt] (13,5.5) -- (19,5.5);
  \filldraw[fill=bkRoute2!15, draw=bkRoute2!80!black, line width=1.1pt] (28,-2.3) -- (35.8,5.5) -- (28,13.3) -- (20.2,5.5) -- cycle;
  \bkqrright{(28,-2.3)}{45}{2.6}
\end{tikzpicture}
\end{minipage}
\end{tsumazuki}

\begin{tsumazuki}{\R{円}{えん}の はばは \R{直径}{ちょっ|けい}}
\R{半径}{はん|けい} 3\,cm の \R{円}{えん}の はしから はしまで\\
\bkbatsu\enspace 3\,cm\hspace{3\zw}\bkmaru\enspace 3 ＋ 3 ＝ 6\,cm\par
\R{円}{えん}の はばは、 \R{中心}{ちゅう|しん}を とおる \R{直径}{ちょっ|けい}。 \R{半径}{はん|けい}の 2つ\R{分}{ぶん}です。
\end{tsumazuki}

\begin{tsumazuki}{\R{直径}{ちょっ|けい}は \R{中心}{ちゅう|しん}を とおる}
\begin{minipage}[c]{0.62\linewidth}
\R{中心}{ちゅう|しん}を とおらない \R{線}{せん}は、 \R{直径}{ちょっ|けい}では ありません。
\end{minipage}\hfill
\begin{minipage}[c]{0.35\linewidth}\centering
\begin{tikzpicture}[x=1mm, y=1mm]
  \draw[bkInk!80, line width=1.1pt] (0,0) circle (9);
  \fill[bkRed] (0,0) circle (0.9);
  \draw[bkRed!80, line width=1pt] (-7.2,5.4) -- (7.2,5.4);
  \draw[bkBlue, line width=1pt] (-9,0) -- (9,0);
  \node at (13,5.4) {\bkbatsu};
  \node at (13,0) {\bkmaru};
\end{tikzpicture}
\end{minipage}
\end{tsumazuki}

\begin{column}{\R{円}{えん}の \R{中心}{ちゅう|しん}を さがそう}
\begin{minipage}[c]{0.7\linewidth}
\R{円}{えん}の \R{紙}{かみ}を \R{半分}{はん|ぶん}に おって、\\もう \R{一度}{いち|ど} \R{半分}{はん|ぶん}に おろう。\\ひらくと、 おりめの \R{交}{まじ}わる\\\R{点}{てん}が \R{中心}{ちゅう|しん}です。
\end{minipage}\hfill
\begin{minipage}[c]{0.27\linewidth}\centering
\begin{tikzpicture}[x=1mm, y=1mm]
  \filldraw[fill=bkYellow!30, draw=bkYellow!70!black, line width=1pt] (0,0) circle (11);
  \draw[bkGray, line width=0.8pt, dash pattern=on 2pt off 1.5pt] (-11,0) -- (11,0) (0,-11) -- (0,11);
  \bkqrright{(0,0)}{0}{2.6}
  \fill[bkRed] (0,0) circle (1);
\end{tikzpicture}
\end{minipage}
\end{column}

\begin{column}{ひもで \R{大}{おお}きな \R{円}{えん}を かこう}
\begin{minipage}[c]{0.6\linewidth}
\R{校庭}{こう|てい}に \R{大}{おお}きな \R{円}{えん}を かく ときは、 ひもを つかいます。
ひとりが ひもの はしを \R{地面}{じ|めん}に おさえて \R{中心}{ちゅう|しん}に なり、 もう ひとりが ひもを ぴんと はったまま ぐるっと \R{歩}{ある}きます。
ひもの \R{長}{なが}さが \R{半径}{はん|けい}。 コンパスと \R{同}{おな}じ しくみです。
\end{minipage}\hfill
\begin{minipage}[c]{0.37\linewidth}\centering
\begin{tikzpicture}[x=1mm, y=1mm]
  \draw[bkGreen!70!black, line width=1pt, dash pattern=on 2.5pt off 1.5pt] (0,0) circle (14);
  \draw[bkStem, line width=1pt] (0,0) -- (40:14);
  \fill[bkRed] (0,0) circle (1.2);
  \fill[bkRoute6] (40:14) circle (1.6);
  \node[font=\sffamily\bfseries\fontsize{7.5}{8}\selectfont, text=bkStem!80!black, fill=white, inner sep=0.5pt] at (40:7) {ひも};
  \draw[bkTheme, line width=0.9pt, -{Stealth[length=3.5pt]}] (60:17) arc (60:110:17);
\end{tikzpicture}
\end{minipage}
\end{column}

\begin{column}{\R{三角形}{さん|かく|けい}は つぶれにくい}
\begin{minipage}[c]{0.56\linewidth}
ストローを 4\R{本}{ほん} つないだ \R{四角形}{し|かく|けい}は、 かどを おすと かたむいて、 \R{形}{かたち}が かわって しまいます。
3\R{本}{ぼん}で つないだ \R{三角形}{さん|かく|けい}は、 \R{辺}{へん}の \R{長}{なが}さが かわらない かぎり、 おしても \R{形}{かたち}が かわりません。
\end{minipage}\hfill
\begin{minipage}[c]{0.42\linewidth}\centering
\begin{tikzpicture}[x=1mm, y=1mm]
  \draw[bkRoute6, line width=1.6pt, line join=round] (0,0) -- (12,0) -- (12,12) -- (0,12) -- cycle;
  \draw[bk arrow, bkGray, line width=0.9pt] (13,6) -- (18,6);
  \draw[bkRoute6, line width=1.6pt, line join=round] (20,0) -- (32,0) -- (38,10.4) -- (26,10.4) -- cycle;
  \node[font=\sffamily\bfseries\fontsize{7.5}{8}\selectfont, text=bkRed] at (29,-3.6) {かたむく};
  \draw[bkRoute4, line width=1.6pt, line join=round] (2,18) -- (16,18) -- (9,30) -- cycle;
  \node[anchor=west, font=\sffamily\bfseries\fontsize{7.5}{8}\selectfont, text=bkRoute4!70!black] at (18,24) {かわらない};
\end{tikzpicture}
\end{minipage}
\end{column}

\clearpage
%% ---------------- p.199 つかわれているところ ----------------
\begin{butsuri}
\R{円}{えん}や \R{三角形}{さん|かく|けい}の ひみつは、 きかいや たてものに いかされて います。
\begin{bkfigure}
\begin{tikzpicture}[x=1mm, y=1mm]
  % はぐるま
  \fill[white, rounded corners=3mm] (0,0) rectangle (39,36);
  \bkqrgear{(13.5,21)}{7.5}{12}{bkRoute6}
  \bkqrgear{(26.8,14)}{5}{8}{bkRoute2}
  \node[font=\sffamily\bfseries\small] at (19.5,4) {はぐるま};
  % マンホール
  \begin{scope}[shift={(42,0)}]
  \fill[white, rounded corners=3mm] (0,0) rectangle (39,36);
  \filldraw[fill=bkGray!45, draw=bkInk!70, line width=1pt] (19.5,20.5) circle (10.5);
  \draw[bkInk!50, line width=0.7pt] (19.5,20.5) circle (8);
  \foreach \a in {0,30,...,150} {\draw[bkInk!40, line width=0.6pt] (19.5,20.5) ++(\a:8) -- ++(\a+180:16);}
  \draw[bkBlue, line width=1.1pt, {Stealth[length=4pt]}-{Stealth[length=4pt]}] (9,20.5) -- (30,20.5);
  \node[font=\sffamily\bfseries\small] at (19.5,4) {マンホールの ふた};
  \end{scope}
  % てっきょう
  \begin{scope}[shift={(84,0)}]
  \fill[white, rounded corners=3mm] (0,0) rectangle (39,36);
  \fill[bkBlue!30] (3,10) rectangle (36,14);
  \draw[bkRoute4!80!black, line width=1.3pt, line join=round] (3,16) -- (36,16)
    (8.5,27) -- (30.5,27) (3,16) -- (8.5,27) -- (14,16) -- (19.5,27) -- (25,16) -- (30.5,27) -- (36,16);
  \node[font=\sffamily\bfseries\small] at (19.5,4) {てっきょう};
  \end{scope}
\end{tikzpicture}
\end{bkfigure}
\begin{itemize}[label={\color{bkPurple}●}, leftmargin=1.4\zw, itemsep=0pt, topsep=0pt]
\item はぐるまや タイヤは \R{円}{えん}。 \R{中心}{ちゅう|しん}が ぶれずに、 なめらかに まわります。
\item まるい ふたは、 どの \R{向}{む}きでも はばが \R{同}{おな}じ。 だから、 あなに おちません。
\item \R{三角形}{さん|かく|けい}の ほねぐみは、 おしても つぶれにくい。 はしや タワーを \R{強}{つよ}く します。
\end{itemize}
\end{butsuri}

\bkfukidashi{delta}{\R{円}{えん}は、 まわる ものの \R{算数}{さん|すう}の はじまり。\\「\R{回転}{かい|てん}と \R{波}{なみ}の \R{路線}{ろ|せん}」に のって いこう！}

\begin{bkfigure}
\begin{tikzpicture}[x=1mm, y=1mm]
  \draw[bkRoute4, line width=3pt, line cap=round] (10,0) -- (122,0);
  \node[anchor=west, inner sep=0pt] at (0,12) {\bkrouteicon[1.4em]{4}\hspace{0.3em}{\sffamily\bfseries\small\color{bkRoute4!80!black}\R{回転}{かい|てん}と \R{波}{なみ}の \R{路線}{ろ|せん}}};
  \foreach \x/\c/\t/\s in {10/bkVol1/{1-15}/{なんしゅう},38/bkVol1/{1-19}/{\R{円}{えん}},66/bkVol1/{1-20}/{\R{角}{かく}},94/bkVol4/{4-8}/{まわる かげ},122/bkVol5/{5-15}/{\R{波}{なみ}の びぶん}} {
    \fill[white, draw=\c, line width=1.6pt] (\x,0) circle (2.6);
    \node[anchor=base, font=\sffamily\bfseries\footnotesize, text=\c!85!black] at (\x,-6.8) {\t};
    \node[anchor=base, font=\sffamily\footnotesize, text=bkInk] at (\x,-13.6) {\s};}
  \node[fill=bkTheme, text=white, rounded corners=1mm, font=\sffamily\bfseries\fontsize{8}{9}\selectfont, inner sep=1.5pt] at (38,5.5) {いま ここ};
\end{tikzpicture}
\end{bkfigure}


\clearpage
%% ---------------- p.200〜201 もんだい ----------------
\bkendproblems

\bkproblemgroup{きほん}{1}

\bkproblem \R{三角形}{さん|かく|けい}と \R{四角形}{し|かく|けい}を えらびましょう。
\begin{bkfigure}
\begin{tikzpicture}[x=1mm, y=1mm]
  \filldraw[fill=white, draw=bkInk!80, line width=1.1pt, line join=round] (0,0) -- (20,2) -- (16,16) -- (2,14) -- cycle;
  \bkqrtag{(-2,19)}{A}
  \filldraw[fill=white, draw=bkInk!80, line width=1.1pt, line join=round] (28,0) -- (46,0) -- (30,16) -- cycle;
  \bkqrtag{(27,19)}{B}
  \draw[bkInk!80, line width=1.1pt, line join=round] (57,3) -- (54,0) -- (72,0) -- (72,15) -- (57,15) -- (57,8);
  \bkqrtag{(53,19)}{C}
  \filldraw[fill=white, draw=bkInk!80, line width=1.1pt, line join=round] (80,0) -- (98,0) .. controls (98,10) and (92,15) .. (86,15) -- cycle;
  \bkqrtag{(79,19)}{D}
  \filldraw[fill=white, draw=bkInk!80, line width=1.1pt, line join=round] (104,6) -- (118,0) -- (114,16) -- cycle;
  \bkqrtag{(105,19)}{E}
\end{tikzpicture}
\end{bkfigure}
\noindent\hspace*{1\zw}\R{三角形}{さん|かく|けい}\enspace\bkkakikomi[22mm]\hspace{1.5\zw}\R{四角形}{し|かく|けい}\enspace\bkkakikomi[22mm]
\bkpans{\R{三角形}{さん|かく|けい} B と E\quad \R{四角形}{し|かく|けい} A}

\bkproblem \R{直角}{ちょっ|かく}は どれ？ \R{三角}{さん|かく}じょうぎで しらべましょう。
\begin{bkfigure}
\begin{tikzpicture}[x=1mm, y=1mm]
  \draw[bkInk!80, line width=1.2pt, line join=round] (0,9) -- (9,0) -- (18,9);
  \bkqrtag{(-1,15)}{A}
  \draw[bkInk!80, line width=1.2pt, line join=round] (44,0) -- (28,0) -- (34,15);
  \bkqrtag{(25,15)}{B}
  \draw[bkInk!80, line width=1.2pt, line join=round] (58,15) -- (58,0) -- (74,0);
  \bkqrtag{(54,15)}{C}
  \draw[bkInk!80, line width=1.2pt, line join=round] (104,0) -- (88,0) -- (80,13);
  \bkqrtag{(84,15)}{D}
\end{tikzpicture}
\end{bkfigure}
\bkkotaeran
\bkpans{A と C}

\bkproblem \R{形}{かたち}の \R{名前}{な|まえ}を かきましょう。
\begin{bkfigure}
\begin{tikzpicture}[x=1mm, y=1mm]
  \bkgrid{(0,0)}{22}{4}{5}
  \filldraw[fill=bkRoute2!20, draw=bkRoute2!80!black, line width=1.2pt] (5,5) rectangle (15,15);
  \bkqrtag{(2.5,17.5)}{A}
  \filldraw[fill=bkRoute6!20, draw=bkRoute6!80!black, line width=1.2pt] (30,5) rectangle (60,15);
  \bkqrtag{(27.5,17.5)}{B}
  \filldraw[fill=bkRoute4!20, draw=bkRoute4!80!black, line width=1.2pt, line join=round] (75,5) -- (100,5) -- (100,15) -- cycle;
  \bkqrtag{(72.5,17.5)}{C}
\end{tikzpicture}
\end{bkfigure}
\noindent A\enspace\bkkakikomi[30mm]\hspace{0.8\zw}B\enspace\bkkakikomi[30mm]\hspace{0.8\zw}C\enspace\bkkakikomi[30mm]
\bkpans{\mbox{A \R{正方形}{せい|ほう|けい}}\quad \mbox{B \R{長方形}{ちょう|ほう|けい}}\quad \mbox{C \R{直角三角形}{ちょっ|かく|さん|かく|けい}}}

\bkproblem \R{半径}{はん|けい} 5\,cm の \R{円}{えん}の \R{直径}{ちょっ|けい}は、 \R{何}{なん}cm ですか。
\bkshikiran
\bkpans{5 × 2 ＝ 10\quad 10\,cm}

\bkproblemgroup{チャレンジ}{1}

\bkproblem この \R{形}{かたち}に、 \R{直角}{ちょっ|かく}は いくつ ありますか。
\begin{bkfigure}
\begin{tikzpicture}[x=1mm, y=1mm]
  \filldraw[fill=bkYellow!20, draw=bkInk!80, line width=1.2pt, line join=round] (0,0) -- (24,0) -- (24,14) -- (12,26) -- (0,14) -- cycle;
\end{tikzpicture}
\end{bkfigure}
\bkkotaeran
\bkpans{3つ（\R{下}{した}の 2つと、 いちばん \R{上}{うえ}の かど）}

\bkproblem \R{直径}{ちょっ|けい} 6\,cm の ボールが 3こ、 つつに ぴったり \R{入}{はい}って います。 つつの \R{長}{なが}さは \R{何}{なん}cm？
\begin{bkfigure}
\begin{tikzpicture}[x=1mm, y=1mm]
  \foreach \x in {6,18,30} {\filldraw[fill=bkYellow!70, draw=bkYellow!60!black, line width=0.8pt] (\x,0) circle (6);
    \draw[white, line width=0.8pt] (\x-4.5,-3.5) .. controls (\x-1,0) and (\x-1,0) .. (\x-4.5,3.5) (\x+4.5,-3.5) .. controls (\x+1,0) and (\x+1,0) .. (\x+4.5,3.5);}
  \draw[bkInk!70, line width=1.1pt] (0,-6.4) -- (36,-6.4) -- (36,6.4) -- (0,6.4) -- cycle;
  \bkbrace{(0,8)}{(36,8)}{？\,cm}{1}
\end{tikzpicture}
\end{bkfigure}
\bkshikiran
\bkpans{6 × 3 ＝ 18\quad 18\,cm}

\begin{furikaeri}
\begin{checklist}
\item \R{三角形}{さん|かく|けい}と \R{四角形}{し|かく|けい}を \R{見}{み}わけられる。\hfill{\footnotesize\sffamily\color{bkTheme}→\pageref{bk:1-19:sankaku}ページ}
\item \R{直角}{ちょっ|かく}や \R{長方形}{ちょう|ほう|けい}・\R{正方形}{せい|ほう|けい}が わかる。\hfill{\footnotesize\sffamily\color{bkTheme}→\pageref{bk:1-19:chokkaku}ページ}
\item \R{円}{えん}の \R{中心}{ちゅう|しん}と \R{半径}{はん|けい}が わかる。\hfill{\footnotesize\sffamily\color{bkTheme}→\pageref{bk:1-19:en}ページ}
\end{checklist}
チェックが つかなかったら、 →の ページに もどろう。
\end{furikaeri}

\bkfukidashi{sora}{うちの \R{中}{なか}の \R{直角}{ちょっ|かく}を さがして みよう。\\\R{三角}{さん|かく}じょうぎを あてると わかるよ！}
\clearpage
